\documentclass{beamer}

\usetheme{Madrid}
\usecolortheme{default}

\usepackage[utf8]{inputenc}
\usepackage[T1]{fontenc}
\usepackage[english]{babel}
\usepackage{listings}
\usepackage{tikz} % Added missing tikz package

\definecolor{ualgblue}{RGB}{0, 74, 130}

\definecolor{codegreen}{rgb}{0,0.6,0}
\definecolor{codegray}{rgb}{0.5,0.5,0.5}
\definecolor{codepurple}{rgb}{0.58,0,0.82}
\definecolor{backcolour}{rgb}{0.95,0.95,0.92}

\lstdefinestyle{mystyle}{
    backgroundcolor=\color{backcolour},   
    commentstyle=\color{codegreen},
    keywordstyle=\color{magenta},
    numberstyle=\tiny\color{codegray},
    stringstyle=\color{codepurple},
    basicstyle=\ttfamily\footnotesize,
    breakatwhitespace=false,         
    breaklines=true,                 
    captionpos=b,                    
    keepspaces=true,                 
    numbers=left,                    
    numbersep=5pt,                  
    showspaces=false,                
    showstringspaces=false,
    showtabs=false,                  
    tabsize=2
}

\title[Minimax Mancuna Agent]{Implementation of a Minimax Agent for Mancuna}
\subtitle{Computer Games - Master in Computer Engineering}
\author[Boavida, Voigt, Schlup]{Alexandre Boavida \and Jan-Erik Voigt \and Philippe Schlup}
\institute[UAlg]{University of Algarve \\ Faculty of Science and Technology}
\date{April 21, 2026}

\begin{document}

\begin{frame}
    \titlepage
\end{frame}

\begin{frame}{Introduction}
    \begin{itemize}
        \item \textbf{Project:} Development of an AI agent for the game of \textit{Mancuna}.
        \item \textbf{Goal:} Implement a competitive player using advanced search techniques.
        \item \textbf{Key Components:}
        \begin{itemize}
            \item Minimax Algorithm
            \item Alpha-Beta Pruning
            \item Transposition Table (Memoization)
            \item Ply-Adjusted Heuristic Evaluation
        \end{itemize}
    \end{itemize}
\end{frame}

\begin{frame}{The Mancuna Game}
    \begin{itemize}
        \item \textbf{Layout:} $2 \times 2$ board (2 pits per player).
        \item \textbf{Tokens:} Typically starts with 2 tokens per pit (8 total).
        \item \textbf{Mechanics:}
        \begin{itemize}
            \item Select a pit, sow tokens counter-clockwise.
            \item Win by emptying the opponent's side on their turn.
            \item Score is the sum of tokens in the player's pits.
        \end{itemize}
    \end{itemize}
\end{frame}

\begin{frame}{Core Algorithm: Minimax \& Pruning}
    \begin{block}{Minimax with Alpha-Beta Pruning}
        \begin{itemize}
            \item Efficiently searches the game tree to find optimal moves.
            \item \textbf{Alpha-Beta Pruning:} Significantly reduces the number of nodes visited by cutting off branches that cannot affect the final decision.
        \end{itemize}
    \end{block}
    \vfill
    \begin{itemize}
        \item \textbf{Max Depth:} Fixed to 12 plies for optimal performance.
        \item Increases in depth did not result in performance gains.
    \end{itemize}
\end{frame}

\begin{frame}{Heuristic Evaluation}
    \begin{itemize}
        \item \textbf{In-Progress State:} Evaluated purely on score differential.
        \begin{itemize}
            \item Calculation: $(\text{Self Score} - \text{Opponent Score}) \times 10$.
        \end{itemize}
        \item \textbf{Terminal States (with Ply Adjustment):}
        \begin{itemize}
            \item \textbf{Tie:} Checked first, strictly returns 0 to prevent the AI from treating ties as losses.
            \item \textbf{Win:} $10000 - \text{ply}$ (Encourages faster wins).
            \item \textbf{Loss:} $-10000 + \text{ply}$ (Encourages longer survival).
        \end{itemize}
        \item The \texttt{ply} parameter strictly represents the distance from the root node, completely solving the traditional Horizon Effect loop trap.
    \end{itemize}
\end{frame}

\begin{frame}{Transposition Table}
    \begin{itemize}
        \item Uses a \texttt{Dictionary} to store previously computed board states.
        \item \textbf{Tuple Hash Key:} \texttt{(board.hash(), turn)}
        \begin{itemize}
            \item Secures the cache by ensuring Max and Min perspectives on the same physical board do not overwrite each other.
        \end{itemize}
        \item \textbf{Stored Data:} Value, Depth, and Flag (Exact, LowerBound, UpperBound).
        \item \textbf{Safety Constraint:} Terminal mate scores ($v \ge 9000$ or $v \le -9000$) are deliberately excluded from caching to prevent depth-related anomalies.
    \end{itemize}
\end{frame}

\begin{frame}[fragile]{Alpha-Beta Pruning Implementation}
\begin{lstlisting}[
    language={[Sharp]C}, 
    basicstyle=\ttfamily\scriptsize, 
    backgroundcolor=\color{gray!10}, 
    frame=single,                    
    rulecolor=\color{gray!40},       
    numbers=left,                    
    numberstyle=\tiny\color{gray},
    xleftmargin=1em                  
]
double v = double.NegativeInfinity;
(List<IBoard> children, _) = board.children();

foreach (var child in children) {
    v = Math.Max(v, MinValue(child, alpha, beta, 
                             remainingDepth - 1, ply + 1));
    // Beta Cutoff
    if (v >= beta) {
        if (v < 9000 && v > -9000) {
            memo[key] = (v, remainingDepth, HashFlag.LowerBound);
        }
        return v;
    }
    alpha = Math.Max(alpha, v);
}
if (v < 9000 && v > -9000)
{
    HashFlag flag = (v > originalAlpha)
    ? HashFlag.Exact : HashFlag.UpperBound;
    memo[key] = (v, remainingDepth, flag);
}

return v;
\end{lstlisting}

\end{frame}



\begin{frame}{Results}
    \begin{itemize}
    \begin{columns}[T]
    % --- Right Column (Bigger for Image) ---
    \begin{column}{1.0\textwidth} % Increased from 0.52
      \centering
    \includegraphics[width=\textwidth]{results.png}
      \vspace{0.2em}
      {\small Obtained results with depth 12}
    \end{column}
  \end{columns}
    \end{itemize}
\end{frame}

\begin{frame}{Experimental Results}
    \begin{itemize}
        \item \textbf{Configuration:} 20,000 matches against an \texttt{Expert} player.
        \item \textbf{Performance Summary:}
        \begin{itemize}
            \item \textbf{Ties:} 20,000 (100\% tie rate).
            \item \textbf{MancunaPlayer0 Total Score:} 100,000.
            \item \textbf{Expert Player Total Score:} 60,000.
        \end{itemize}
        \item \textbf{Analysis:}
        \begin{itemize}
            \item The 100\% tie rate indicates that both agents are playing optimally given the small state space (8 tokens, 4 pits).
            \item Our agent achieved a higher cumulative score while maintaining a perfect defensive record.
            \item Consistent performance across 250,000+ moves.
        \end{itemize}
    \end{itemize}
\end{frame}

\begin{frame}{Future Work \& Possible Improvements}
    \begin{itemize}
        \item \textbf{Iterative Deepening:}
        \begin{itemize}
            \item Transitioning from a fixed-depth search to a time-bounded search.
            \item Acts as an "anytime algorithm" to guarantee an optimal move within strict time limits.
        \end{itemize}
        \item \textbf{Deep Move Ordering:}
        \begin{itemize}
            \item Sorting child nodes to evaluate the strongest moves first, driven by the Transposition Table's previous findings.
            \item Maximizes Alpha-Beta pruning efficiency, vastly reducing the number of nodes evaluated.
        \end{itemize}
        \item \textbf{Advanced Search Enhancements:}
        \begin{itemize}
            \item \textbf{Aspiration Windows:} Narrowing the initial $\alpha$ and $\beta$ bounds to trigger fail-high/fail-low cutoffs faster.
            \item \textbf{Endgame Tablebases:} Precomputing mathematically perfect moves for low-token scenarios to bypass search entirely.
        \end{itemize}
    \end{itemize}
\end{frame}

\begin{frame}{Conclusion}
    \begin{itemize}
        \item The \texttt{PlayerMancuna} agent provides a robust implementation of classic AI game-playing techniques.
        \item Alpha-Beta pruning combined with a Transposition Table allows for highly efficient deep look-ahead up to 12 plies.
        \item Ply-adjusted terminal heuristics ensure logical handling of forced game endings.
    \end{itemize}
    \vfill
    \centering
    
\end{frame}

% ============================================================
\section*{Q\&A}
% ============================================================

\begin{frame}[plain]
  \begin{tikzpicture}[remember picture, overlay]
    \fill[ualgblue] (current page.south west) rectangle (current page.north east);
  \end{tikzpicture}

  \centering
  \vspace{2.8cm}
  {\Huge\color{white}\textbf{Any Questions?}}\\[1.2em]
  {\large\color{white!85}Thank you for your attention}\\[2.2em]

  % \begin{beamercolorbox}[sep=0.6em, center, wd=0.65\textwidth,
  %     rounded=true]{palette primary}
  %   \small
    
  % \end{beamercolorbox}
\end{frame}

\end{document}
